\section{A ternary five-subset network}
\label{sec:ternary-five-subsets}

This section changes the finite scalar alphabet for the bit interchange
network. It retains the rational-frame contract and the compact-control
movement proof. The complex network and its guard constants are unchanged.
The result remains conditional on the pinned upstream manuscript and the
previously identified compact-control and Gaussian extensions.

Zhihao Chen's earlier
\href{https://github.com/CrocSwap/integer-mult-bounds/pull/7}{PR~\#7}
introduces the same ternary five-subset motif, with a different producer and
a stronger claimed witness. The circuit below is a separate implementation;
no priority claim for the motif or verification of that pending PR is made.

\subsection{The two fields and the side map}
Let $h=29$, let $\mathcal T=\binom{[h]}5$, and put
$n=\binom h5=118755$, $r=\binom h2=406$. Over $\mathbb F_3$ let $B$ be
the incidence matrix of five-subsets against pairs. Then
\[
 C=BB^{\mathsf T},\qquad C_{ST}=\binom{|S\cap T|}{2}\pmod3.
\]
For $t=0,1,2,3,4,5$ the binomial residues are $0,0,1,0,0,1$.
Consequently $C=I+A_2$, where $A_2$ is the adjacency matrix for
intersection exactly two. The side map is $-A_2$; thus the central
and side maps add to the identity. Only the displayed factor size $r$
is needed, not a claim that it is a minimal factor.

The rational label space is $F=\mathbb Q^h$ with form
\[
 H=I-\frac2{25}J.
\]
It is nondegenerate, since $1-2h/25=-33/25\ne0$. Write $t_T$ for the
indicator of $T$. Then
\[
 t_S^{\mathsf T}Ht_T=|S\cap T|-2,\qquad
 t_T^{\mathsf T}Ht_T=3.
\]
Let $P_T$ be projection onto $\langle t_T\rangle$ in this form.
Intersection-two neighbors have mutually annihilating projections.
The scalar characteristic and rational label field are different;
the address prime in the matrix-shear procedure is a third, independently
chosen fixed prime avoiding its rational denominators.

\subsection{An aligned side sum circuit}
For each pair $J\subset[h]$, take all source variables $x_{J\cup S}$
with $|S|=3$ and $S\cap J=\varnothing$. For each triple $T$ disjoint
from $J$ compute
\[
 z_{J,T}=\sum_{S\cap(J\cup T)=\varnothing,\ |S|=3}x_{J\cup S}.
\]
Inject $-z_{J,T}$ into target $J\cup T$. An ordered intersection-two
pair has the unique common pair $J$, so these partial outputs implement
$-A_2$ exactly. There are $q=\binom h2\binom{h-2}3=10n$ output uses.

Use the deterministic disjoint-triple sum DAG of the tensor construction
on all $h$ points. In group $J$, delete formal source variables meeting
$J$ and retain only output triples disjoint from $J$. A zero in this
substitution denotes an absent formal input, not a cleared physical
scratch register. Every retained addition still has disjoint source
supports. The fixed global partition order aligns intermediate sums
between different common-pair groups.

For completeness, the tensor sum recipe is a finite recursion on ordered
ground sets. Let $Z(n,a,b)$ be its uninterned addition count for summing
$a$-subset inputs over all $b$-subset disjointness queries. If $a+b>n$,
the outputs are empty sums; if $a=0$, copy the sole input. Both have
zero cost. If $b=0$, sum the $\binom na$ inputs at cost $\binom na-1$.
Otherwise start with separate query sums, costing
$\binom nb(\binom{n-b}a-1)$. For each $1\le l\le\lfloor n/2\rfloor$,
put $r_0=n-l$ and split the ground set into its first $l$ and last $r_0$
points. For every feasible input split $i$ and query split $j$, contract
one factor and then the other, at cost
\[
 \min\left\{
 \binom{r_0}{a-i}Z(l,i,j)+\binom lj Z(r_0,a-i,b-j),\quad
 \binom li Z(r_0,a-i,b-j)+\binom{r_0}{b-j}Z(l,i,j)
 \right\}.
\]
Feasibility requires the subset sizes to fit their factors and
$i+j\le l$, $a-i+b-j\le r_0$. For each $j$, add the resulting terms
for each full query: if there are $t_j>0$ feasible $i$, this costs
$\binom lj\binom{r_0}{b-j}(t_j-1)$. Choose a split only if its total
strictly improves the current best, testing $l$ in increasing order.
For the displayed minimum prefer its first contraction order on a tie.
Use lexicographic subset order and increasing split indices; sum lists
by recursively splitting them into two halves. Intern every equal formal
support at each addition, retaining its first decomposition. The initial
template is the resulting $(n,a,b)=(h,3,3)$ DAG. This describes a
cancellation-free integer coefficient computation: disjointness factors
across the two ground sets, and the input split partitions the sources.

Within a group, intern equal full formal support sets. Across groups,
let $I$ be the intersection of all full five-subset source labels of a
nontrivial sum. A sum shared by distinct groups has $|I|\ge3$.
If $|I|=3$, its exact signature is $I$ and the set of residual pairs
$S\setminus I$; if $|I|=4$, use $I$ and its residual points. These
signatures reconstruct the full set of five-subset variables, so there
are no hash collisions at the mathematical level. A nontrivial sum
cannot have $|I|=5$. Sums with $|I|=2$ belong only to that common-pair
group. Equal singleton inputs are identified by their five-subset labels.
Retain the first decomposition of each sum and prune unused nodes after
all groups are processed.

The exact integer DAG-array construction at $h=29$ gives
\[
 c=19593239,\qquad q=1187550,\qquad S=c+q=20780789.
\]
Here $c$ counts active additions after pruning, not allocated nodes.
The implementation independently expands every support and output row
at $h=8,9$, checks the signature identifications there, and regenerates
the full integer DAG counts at $h=29$. The general exact-signature
argument above justifies the large sharing operation; the small tests
are implementation controls, not an extrapolated proof.

\subsection{Restoration and rational side frames}
A cancellation-free addition DAG with $c$ additions and $q$ output uses
has an invertible gather-and-fanout embedding on $S=c+q$ roles: each
binary addition identifies one input pivot with an output use, while
all other uses receive their own role. No initial side value is assumed
zero. Let its invertible scalar operation be $L$, its data copy be $V$,
and its signed output injection be $J$ (coefficient $-1$ at every partial
output). With central gather $G=B^{\mathsf T}$ and scatter $R=B$, use
\[
 L,\ -J,\ L^{-1},\ -R,\ V,\ G,\ R,\ L,\ J,\ L^{-1},\ -G,\ -V.
\]
On arbitrary side input $z$, the two injections contribute
$-JLz+JL(z+Vx)=JLVx=-A_2x$. The inverse mixers and final negative copy
restore $z$. Central input cancels in the same way, leaving $RGx=Cx$.
Thus $y\gets y+x$ and every auxiliary input is restored. The actual
reverse inverse schedule gives $y\gets y-x$.

Every nonzero vector $u$ in the span of five-subset indicators containing
one fixed pair $\{a,b\}$ satisfies
\[
 u_a=u_b,\quad \sum_i u_i=5u_a,\quad
 u^{\mathsf T}Hu=\sum_{i\notin\{a,b\}}u_i^2>0.
\]
If the last sum vanished, the first two identities would give
$2u_a=5u_a$, hence $u=0$. Therefore every source span appearing in an
original group is nondegenerate. Interning an equal formal sum does not
change its span, so this remains true after cross-group sharing.

For a DAG node $v$ choose the orthogonal projection $E_v$ onto its source
span. These projections increase along forward edges. Every ancestor
source is orthogonal to every reachable target, so the final projection
at a partial output lies below $I-P_T$ for its target $T$. For the
reverse mixer use $I-E_v$. Orthogonal complementation reverses the
inclusions and gives increasing reverse edges, including the target
lines at entry and source complements at exit. Rational nondegeneracy
suffices; no characteristic-two orthonormal-residual assumption is used.

\subsection{Complete tensor circuit and rank ledger}
Index the data banks by $\mathcal T^3$, with $N=n^3$ and $m=h^3$.
Use the first-coordinate shears, inverse second-coordinate shears with
logical banks exchanged, and third-coordinate shears. The scalar action
is $(X,Y)\mapsto(-Y,X)$. Negate the physical $X$ bank at its final identity
frame to obtain the unsigned bank exchange. This fixed pointwise gate
has zero added edge rank and leaves all restored auxiliaries unchanged.

For a stage write $A=B_0\mathbin\perp P$, where $P$ is its preceding
label tensor line, and suppress the future tensor line $Q$. Put
\[
 D_0=B_0\otimes I,\quad D_1=A\otimes I,\quad
 D(E)=D_0+P\otimes E.
\]
The common frames of the forward twelve operations are
\[
 D_0,D_0,D_0,D_0,D(P_T),D_1,D_0,D(E_v),D(I-P_T),D_1,D_1,D_1.
\]
For the reversed inverse invocation they are
\[
 D_0,D_0,D_0,D(P_T),D(I-E_v),D_1,D_0,D(I-P_T),D_1,D_1,D_1,D_1.
\]
Copy and injection gates are grouped by their data role, including all
partial outputs for that target. Central gates retain every data and
central incidence, including zero factor coefficients. All frames are
tensored with $Q$. These are the same tensor boundaries as the general
three-stage compiler. All side changes increase; only the central return
decreases, by $h$ on each of the $r$ central roles in each invocation.
Thus the total decreasing dimension is
\[
 L_{\mathrm{dec}}=3n^2rh.
\]
The signed dimension telescoping gives $Wm-2N+2L_{\mathrm{dec}}$.
Replace the physical $X$ source projections $P_a$ by $-P_a$; because
$2\ne0$ in the rational frame field, this adds exactly $N$ to the rank
sum. Source frames are now $-P_a,0$, sink frames $I,I-P_a$, and each
restored auxiliary has source $0$, sink $I$. The unsigned scalar exchange
therefore satisfies every endpoint difference $I_m$.

The bipartite double cover of the intersection-two graph is regular of
positive degree $10\binom{h-5}3$. Hall's condition follows by edge
counting, and a fixed deterministic matching algorithm supplies an
orthogonal label permutation $\pi$. Share stage-one invocation $(A,B)$
with stage-three invocation $(B,\pi(A))$, keeping local auxiliary indices.
The boundary projections are
\[
 E=I\otimes P_A\otimes P_B,\qquad
 H_0=(I-P_B\otimes P_{\pi(A)})\otimes I.
\]
They satisfy $EH_0=H_0E=E$. The twelve-operation schedule already has
initial frame $D_0$ on every auxiliary, so the join replaces two boundary
segments by one and saves exactly $m$ ranks per identified role.
Consequently
\[
 W=2n^3+2n^2(S+r),\quad
 \Delta=Wm-s=n^3-6n^2rh.
\]
Numerically,
\[
\begin{split}
 m&=24389,\qquad N=1674772079218875,\\
 W&=589493540769997500,\\
 L_{\mathrm{dec}}&=498137336383050,\\
 \Delta&=678497406452775,\\
 s&=14377157287342062574725.
\end{split}
\]
Writing $\eta=\Delta/(Wm)$, the exact logarithm enclosure verifies
\[
 \eta>\frac{467}{10^{11}}\log m.
\]
Since $1-\eta<e^{-\eta}$, this proves
$s/W<m^{1-a_{\rm b}}$ with $a_{\rm b}=467/10^{11}$.
The finite-alphabet transfer therefore supplies the bit interchange
exponent $\tau=1-467/10^{11}$.

\subsection{Conditional assembly witness}
Retain the separately integrated complex network with
$a_{\rm c}=5/10^9$, its arity $17576$, and its existing node charge.
Its stopped-depth constants remain $\beta=1/100$, $\zeta=1/1000$ and
$C_1=4961/1000$. The bit alphabet extension is exact and introduces no
complex rounding operations. Its fixed symbol width and circuit size
alter constants only, not the complex precision recurrence.

Take
\[
 \epsilon=\frac{1999}{10000},\quad c=1,\quad\delta=10^{-6},\quad
 \lambda=1-\frac{4669}{10^{12}},\quad
 \lambda'=1-\frac{4668}{10^{12}},\quad\kappa=2^{-30}.
\]
Since $\sigma<\tau$, the compact-control internal exponent is
$\chi=\tau$. The reserved-axis preprocessing exponent is zero. Exact
rational arithmetic checks every compact-control, Gaussian, leaf,
prefix, dimension, guard and assembly constraint, with these two proved
substitutions for the old packed overhead. All seven final margins exceed
$\kappa$; the minimum is
\[
 g_3=\frac{2332833}{2500000000000000},\qquad
 g_3-2^{-30}=\frac{296652863}{163840000000000000000}>0.
\]
Subject to the retained upstream and written extension hypotheses, this
therefore supports $T(n)=O(n(\log n)^{1-2^{-30}})$. It is an asymptotic
exponent comparison, not a practical runtime claim or independent
verification of the upstream multiplication theorem.
